\documentclass[a4paper]{article}
\usepackage{amsfonts}
\usepackage{amsmath}
\usepackage{amssymb}
\usepackage{graphicx}     

\begin{document}
  
\noindent \large Auteur: Sadia Nazary        \\
\noindent \large Studierichting: Informatica \\
\noindent \large Oefeningenreeks 2B nr. 5.3 \\

\medskip

\normalsize

$ \operatorname{Bgtan}(x) + \operatorname{Bgtan}\left( \dfrac{1}{x} \right) $ \\

Stel $ x = \tan(\alpha) $ met $ \alpha \in \left]0, \dfrac{\pi}{2}\right[ $ \\

$\Rightarrow \dfrac{1}{x} = \cot(\alpha) = \tan\left( \dfrac{\pi}{2} - \alpha \right) $ \\

Dus $ \operatorname{Bgtan}\left( \dfrac{1}{x} \right) = \dfrac{\pi}{2} - \alpha $ \\

$ \Rightarrow \operatorname{Bgtan}(x) + \operatorname{Bgtan}\left( \dfrac{1}{x} \right) = \alpha + \left( \dfrac{\pi}{2} - \alpha \right) = \dfrac{\pi}{2} $ \\

\textbf{Conclusie:} \\

$$ \forall x > 0 : \operatorname{Bgtan}(x) + \operatorname{Bgtan}\left( \dfrac{1}{x} \right) = \dfrac{\pi}{2} $$\\

\begin{thebibliography}{999}
%\bibitem{a} Referentie
\end{thebibliography}

\end{document}
